% Section du chapitre 4 (ex-annexe D)
% Annexe issue de l'ancien chapitre 3 (restructuration en 3 chapitres)
% ============================================================================
% CHAPITRE 3 — SECTION 3.4 : VALIDATION RÉTROSPECTIVE
% Fichier : partie3/donnees_validation.tex
% ============================================================================

\section{VALIDATION R\'ETROSPECTIVE}
\label{sec:validation_retrospective}
% ============================================================================
\begin{sloppypar}

La section pr\'ec\'edente a sp\'ecifi\'e \emph{comment} mesurer les
param\`etres institutionnels. La pr\'esente section \'etablit \emph{quelles
valeurs} ces param\`etres prennent, par estimation r\'etrospective sur les
trois syst\`emes de r\'ef\'erence classiques et les standards
post-quantiques. Aucune mesure nouvelle n'est introduite, de sorte que les chiffres qui
suivent sont des calibrations illustratives, d\'eriv\'ees des sp\'ecifications
publi\'ees et mises en regard des produits CLO de la
section~\ref{sec:derivation_opposabilite}.
\end{sloppypar}

\begin{boiteArbitre}
L'arbitre change de posture, car apr\`es un plan, il assiste \`a une
confrontation. Ce qui avait \'et\'e admis plus haut ($\Delta \in [35, 73]$
bits, $\alpha \in [0{,}8, 1{,}5]$, $\beta = O(1)$) revient maintenant sous
forme de valeurs calibr\'ees, syst\`eme par syst\`eme. Il s'agit de savoir
si les ordres de grandeur retenus se tiennent entre eux, sans v\'erifier le
th\'eor\`eme, dont ces syst\`emes ne sont pas des instances.
\end{boiteArbitre}

% -----------------------------------------------------------------------
\subsection{Estimation d\'etaill\'ee des param\`etres institutionnels classiques}
\label{sec:estimation_params_classiques}
% -----------------------------------------------------------------------
\begin{sloppypar}

La section~\ref{sec:derivation_opposabilite} a d\'eriv\'e les produits CLO des
trois syst\`emes de r\'ef\'erence par la voie descendante, en partant des
espaces de verdicts. La pr\'esente section emprunte la voie oppos\'ee et
d\'ecompose les param\`etres institutionnels $(\Delta, \alpha, \beta)$ qui
param\`etrent les bornes du chapitre~\ref{chap:cadre_formel}. Les deux voies ne
sont pas ind\'ependantes (section~\ref{sec:synthese_derivations})~; leur mise en
regard est une calibration.

\subsubsection{Composantes g\'en\'eriques de la capacit\'e d'extraction}
\label{subsec:composantes_generiques}

La capacit\'e d'extraction institutionnelle $\Delta$ agr\`ege six familles de
ressources, dont le tableau~\ref{tab:composantes_delta} donne les plages
observ\'ees dans la litt\'erature cognitive et organisationnelle.

\begin{table}[htbp]
\centering
\caption{Familles de composantes de la capacit\'e d'extraction institutionnelle
         $\Delta$ et plages associ\'ees}
\label{tab:composantes_delta}
\begin{tabular}{@{}lcc@{}}
\toprule
\textbf{Composante} & \textbf{Plage} & \textbf{Base} \\
\midrule
Pr\'ec\'edents l\'egaux et jurisprudence & 10--15 bits & Cat\'egories de cas document\'es \\
Formation d'expert et connaissances & 10--20 bits & Domaine de sp\'ecialisation \\
R\`egles proc\'edurales et m\'emoire institutionnelle & 5--10 bits & Protocoles \'ecrits \\
D\'elib\'eration collaborative & 0--15 bits & Amplification collective \\
Capacit\'e temporelle de d\'elib\'eration & 5--10 bits & Temps allou\'e \\
Ressources externes et consultation & 5--15 bits & Experts, r\'ef\'erences \\
\bottomrule
\end{tabular}
\end{table}

\noindent\textbf{Lecture du tableau.} Ces plages ne sont pas cumulatives, car une
institution donn\'ee ne mobilise qu'un sous-ensemble de ces familles. Certaines se
recouvrent d'ailleurs partiellement, puisque la consultation d'experts externes
et la formation interne alimentent la m\^eme ressource cognitive. La
sommation na\"{\i}ve des bornes sup\'erieures produirait $85$ bits, valeur
sup\'erieure \`a toutes les estimations retenues dans cette th\`ese. Le
tableau~\ref{tab:decomposition_delta} proc\`ede donc syst\`eme par syst\`eme,
en ne retenant que les composantes effectivement mobilis\'ees dans chaque
contexte institutionnel.

\subsubsection{D\'ecomposition de $\Delta$ par syst\`eme}
\label{subsec:decomposition_delta}

La capacit\'e d'extraction institutionnelle $\Delta$ est d\'ecompos\'ee, \`a titre
illustratif, en composantes cognitives et organisationnelles. Les totaux
($15+10+10+0 = 35$, $15+20+15+6 = 56$, $25+20+15+13 = 73$) sont ceux de la table
de r\'ef\'erence (tableau~\ref{tab:parametres_reference})~; la d\'ecomposition du
tableau~\ref{tab:validation_cro_classique} du chapitre~\ref{chap:cadre_formel} y renvoie.

\begin{table}[htbp]
\centering
\caption{D\'ecomposition d\'etaill\'ee de $\Delta$ (bits) par syst\`eme classique}
\label{tab:decomposition_delta}
\begin{tabular}{@{}lcccl@{}}
\toprule
\textbf{Composante de $\Delta$} & \textbf{ECDSA} & \textbf{Helios} &
\textbf{zk-SNARK} & \textbf{Justification} \\
\midrule
Pr\'ec\'edent l\'egal / jurisprudence
  & 15 & 15 & 25
  & Doctrine \'etablie \\
Expertise technique / formation
  & 10 & 20 & 20
  & Formation sp\'ecialis\'ee \\
Proc\'edures / cadres institutionnels
  & 10 & 15 & 15
  & Standards, protocoles \\
D\'elib\'eration collaborative
  & 0 & 6 & 13
  & Comit\'es multi-parties \\
\midrule
\textbf{Total $\Delta$} & \textbf{35} & \textbf{56} & \textbf{73} & \\
\midrule
\textbf{R\'egime (Rem.~\ref{rem:activite_delta})} & rare & rare & riche & \\
\bottomrule
\end{tabular}
\end{table}

\subsubsection{Coh\'erence avec la rationalit\'e born\'ee}
\label{subsec:coherence_rationalite}

Les valeurs $\Delta \in [35, 73]$ bits correspondent \`a des facteurs
d'amplification de $1{,}8$ \`a $3{,}7$ par rapport \`a une capacit\'e
individuelle de r\'ef\'erence fix\'ee \`a $\approx 20$ bits. Ce rep\`ere est une
analogie inspir\'ee des travaux sur les limites de la m\'emoire de
travail~\cite{Miller56}. Miller ne l'a pas mesur\'e~; la d\'ecomposition
ci-dessous est donc heuristique.

\begin{equation}
\text{Amplification} = 1{,}5_{\text{doc}} \times 1{,}5_{\text{collab}} \times 1{,}3_{\text{sp\'ec}} \approx 2{,}9\times
\end{equation}

\subsubsection{Param\`etres $\alpha$ et $\beta$}
\label{subsec:parametres_alpha_beta}

\begin{table}[htbp]
\centering
\caption{Param\`etres $\alpha$ et $\beta$ par syst\`eme classique}
\label{tab:parametres_alpha_beta}
\begin{tabular}{@{}lcccl@{}}
\toprule
\textbf{Syst\`eme} & $\alpha$ & $\beta$ & \textbf{Variante SOTP} &
\textbf{Justification} \\
\midrule
ECDSA & 1{,}0 & 1{,}0 & SOTP-Low
  & Validit\'e binaire, nonce uniforme \\
Helios & 1{,}2 & 1{,}4 & SOTP-Base
  & Audit d\'etaill\'e, irr\'egularit\'e ZK \\
zk-SNARKs & 1{,}5 & 1{,}5 & SOTP-High
  & Conformit\'e gradu\'ee, circuit optimis\'e \\
\bottomrule
\end{tabular}
\end{table}
\end{sloppypar}

% -----------------------------------------------------------------------
\subsection{Estimation des param\`etres quantiques}
\label{sec:estimation_params_quantiques}
% -----------------------------------------------------------------------
\begin{sloppypar}

\subsubsection{D\'ecomposition de $\Delta_q$ par syst\`eme PQC}
\label{subsec:decomposition_delta_q}

La capacit\'e d'extraction quantique $\Delta_q$ int\`egre les modifications
induites par la transition quantique (section~\ref{sec:cro_pq}).

\begin{table}[htbp]
\centering
\caption{D\'ecomposition d\'etaill\'ee de $\Delta_q$ (bits) par syst\`eme PQC}
\label{tab:decomposition_delta_q}
\footnotesize
\begin{tabular}{@{}lcccl@{}}
\toprule
\textbf{Composante} & \textbf{ML-DSA-65} & \textbf{ML-KEM-1024}$^\dagger$ &
\textbf{SLH-DSA} & \textbf{Justification} \\
\midrule
Capacit\'e classique $\Delta$ & 35 & 45 & 30 & Contexte institutionnel \\
P\'enalit\'e d'adaptation $(-\delta_{\text{degr}})$ & $-7$ & $-6$ & $-5$ & Nouveau paradigme \\
Am\'elioration quantique $(+\delta_{\text{enh}})$ & $+0$ & $+0$ & $+0$ & Pas encore exploitable \\
\midrule
\textbf{Total $\Delta_q$} & \textbf{28} & \textbf{39} & \textbf{25} & \\
\midrule
\textbf{Plage NIST PQC} & \multicolumn{3}{c}{$\Delta_q \in [25, 67]$ bits} & Hypoth\`ese \\
\bottomrule
\end{tabular}

\smallskip
{\footnotesize $^\dagger$ML-KEM est un m\'ecanisme d'encapsulation de cl\'e, qui ne produit
pas de preuve. Il est donc hors du champ du Trilemme et n'est conserv\'e que pour la
tra\c{c}abilit\'e avec les valeurs publi\'ees (fichier \texttt{06\_calcul/data/calibration\_pq.csv}).\par}
\end{table}

\noindent\textbf{Justification de la p\'enalit\'e d'adaptation.} La transition
vers les standards NIST PQC impose une courbe d'apprentissage aux institutions.
La p\'enalit\'e $\delta_{\text{degr}}$ refl\`ete~:
\begin{itemize}
\item Nouveau format de cl\'es (structur\'e diff\'eremment de RSA/ECDSA)
\item Nouveaux algorithmes de v\'erification (lattice-based vs discrete log)
\item Manque de pr\'ec\'edents l\'egaux pour les signatures post-quantiques
\item Documentation technique encore en maturation
\end{itemize}

\subsubsection{Estimation de $\eta_q$ par niveau de bruit quantique}
\label{subsec:estimation_eta_q}

Pour les signatures post-quantiques, qui sont des cha\^{\i}nes de bits
classiques, $\eta_q = 0$. Le tableau~\ref{tab:estimation_eta_q} d\'ecrit le
\emph{sc\'enario hypoth\'etique} d'un support ou d'un v\'erificateur quantique~;
les valeurs sont des ordres de grandeur a priori.

\begin{table}[htbp]
\centering
\caption{Estimation de $\eta_q$ par type de canal quantique}
\label{tab:estimation_eta_q}
\begin{tabular}{@{}lcccc@{}}
\toprule
\textbf{Canal / Environnement} & \textbf{Param\`etre} & \textbf{Fid\'elit\'e $F$} &
$\eta_q$ \textbf{estim\'e} & \textbf{Source} \\
\midrule
Syst\`eme corrig\'e (id\'eal) & $p \leq 10^{-4}$ & $\geq 0{,}9999$ & $0{,}005$--$0{,}01$ & Th\'eorique \\
NISQ actuel (IBM, IonQ) & $p \approx 10^{-3}$ & $\approx 0{,}99$ & $0{,}05$--$0{,}10$ & Mesures 2024--2025 \\
NISQ bruit\'e & $p \approx 10^{-2}$ & $\approx 0{,}95$ & $0{,}10$--$0{,}20$ & Sc\'enario pessimiste \\
Simulation hybride & Variable & Variable & $0{,}02$--$0{,}15$ & Param\'etrable \\
\midrule
\textbf{Valeur de sc\'enario} & \multicolumn{3}{c}{$\eta_q \approx 0{,}08$ (m\'ediane NISQ)} & \cite{MINKA2026PQCRO} \\
\bottomrule
\end{tabular}
\end{table}

\noindent\textbf{Relation heuristique.}
\begin{equation}
\eta_q \lesssim 2\sqrt{1 - F_{\mathrm{avg}}}
\label{eq:relation_eta_q_fidelite}
\end{equation}
Cette relation n'a pas valeur de th\'eor\`eme. L'in\'egalit\'e de Fuchs--van de Graaf
borne la distance de trace par $\sqrt{1-F}$ tandis que la continuit\'e de l'information
mutuelle conditionnelle (Alicki--Fannes--Winter) convertit cette distance en perte
d'information avec un facteur qui d\'epend de la dimension du syst\`eme au lieu de
$H(C)$. Elle fixe un ordre de grandeur pour le sc\'enario.

\subsubsection{Param\`etres $\alpha_q$ et $\beta_q$ pour PQC}
\label{subsec:parametres_alpha_q_beta_q}

\begin{table}[htbp]
\centering
\caption{Param\`etres post-quantiques par syst\`eme NIST PQC ($H(\mathcal{C})$ \`a l'\'echelle de l'esp\`ece)}
\label{tab:parametres_alpha_q_beta_q}
\begin{tabular}{@{}lcccl@{}}
\toprule
\textbf{Syst\`eme} & $\alpha_q$ & $\beta_q$ & $H(\mathcal{C})$ &
\textbf{Justification} \\
\midrule
ML-DSA-65 & 1{,}0 & 1{,}2 & 10{,}5 bits & Signature (FIPS~204) \\
ML-KEM-1024$^\dagger$ & 1{,}1 & 1{,}3 & 12{,}0 bits & KEM, hors champ \\
SLH-DSA & 0{,}9 & 1{,}1 & 9{,}0 bits & Signature (FIPS~205) \\
\bottomrule
\end{tabular}
\end{table}
\end{sloppypar}

% -----------------------------------------------------------------------
\subsection{Comparaison des produits CLO classiques et quantiques}
\label{sec:comparaison_produits_cro}
% -----------------------------------------------------------------------
\begin{sloppypar}

Le tableau~\ref{tab:comparaison_produits_cro} rassemble les produits CLO
illustratifs des syst\`emes \'etudi\'es. Pour les trois syst\`emes classiques,
la colonne \guillemotleft~Produit CLO~\guillemotright{} reprend les valeurs
d\'eriv\'ees \`a la section~\ref{sec:derivation_opposabilite} (avec $\Relver$).
Pour les standards post-quantiques, o\`u aucune d\'erivation descendante n'est
disponible faute de contexte institutionnel document\'e, les produits sont
obtenus par transposition du rapport $\Delta_q/\lambda$ major\'e d'une correction
d'adaptation, si bien qu'ils constituent des hypoth\`eses de travail.

% -----------------------------------------------------------------------
\subsection{Tableau comparatif principal}
\label{sec:tableau_comparatif_principal}
% -----------------------------------------------------------------------

\begin{table}[htbp]
\centering
\caption{Produits CLO illustratifs, classiques et post-quantiques ($\lambda = 128$)}
\label{tab:comparaison_produits_cro}
\begin{tabular}{@{}llcccc@{}}
\toprule
\textbf{Type} & \textbf{Syst\`eme} & $\Delta/\Delta_q$ & $\eta_q$ (sc\'enario) &
\textbf{Produit CLO} & \textbf{\'Ecart vs classique} \\
\midrule
\multicolumn{6}{c}{\textit{Syst\`emes classiques (calibration)}} \\
\midrule
Classique & ECDSA & 35 & --- & 0{,}273 & --- \\
Classique & Helios & 56 & --- & 0{,}437 & --- \\
Classique & zk-SNARKs & 73 & --- & 0{,}572 & --- \\
\midrule
\multicolumn{6}{c}{\textit{Syst\`emes post-quantiques (hypoth\`ese de travail)}} \\
\midrule
PQC & ML-DSA-65 & 28 & 0{,}08 & 0{,}35 & --- \\
PQC & ML-KEM-1024$^\dagger$ & 39 & 0{,}08 & 0{,}42 & --- \\
PQC & SLH-DSA & 25 & 0{,}08 & 0{,}30 & --- \\
\midrule
\multicolumn{6}{c}{\textit{Comparaison dans des usages voisins}} \\
\midrule
Classique $\to$ PQC & ECDSA $\to$ ML-DSA & 35 $\to$ 28 & 0 $\to$ 0{,}08 &
  $0{,}273 \to 0{,}35$ & $+28\%$ \\
Classique $\to$ PQC & RSA $\to$ SLH-DSA & 40 $\to$ 25 & 0 $\to$ 0{,}08 &
  $0{,}35 \to 0{,}30$ & $-14\%$ \\
\bottomrule
\end{tabular}

\smallskip
{\footnotesize $^\dagger$Hors champ (m\'ecanisme d'encapsulation de cl\'e).\par}
\end{table}

\noindent\textbf{Observation.} La p\'enalit\'e d'adaptation ($\Delta_q < \Delta$)
abaisse la borne. Le terme $\eta_q H(C)$, qui n'intervient que dans le sc\'enario
d'un support quantique, la rel\`eve l\'eg\`erement. Le produit
ML-DSA augmente pourtant par rapport \`a ECDSA alors que $\Delta_q$ diminue, ce qui tient
au fait que la transposition utilis\'ee pour ces produits n'est pas d\'eriv\'ee de la borne.

\subsubsection{Analyse de la contribution du terme $\eta_q \cdot H(C)$}

La majoration affaiblie $(\Delta_q + \eta_q \cdot H(C))/\lambda$ agr\`ege deux
termes d'ordres de grandeur tr\`es diff\'erents. Le
tableau~\ref{tab:decomposition_borne_pq} en donne la d\'ecomposition et la
confronte au produit illustratif du tableau~\ref{tab:comparaison_produits_cro}.

\begin{table}[htbp]
\centering
\caption{D\'ecomposition de la majoration CRO$^q$ et confrontation au produit
         illustratif ($\lambda = 128$, sc\'enario $\eta_q > 0$)}
\label{tab:decomposition_borne_pq}
\small
\begin{tabular}{@{}lcccccc@{}}
\toprule
\textbf{Syst\`eme} & $\Delta_q/\lambda$ & $\eta_q \cdot H(C)/\lambda$ &
\textbf{Borne} & \textbf{Part $\eta_q$} & \textbf{Produit} & \textbf{D\'epass\'ee~?} \\
\midrule
ML-DSA-65 & $0{,}219$ & $0{,}006$ & $0{,}225$ & $2{,}7\%$ & $0{,}35$ & oui \\
ML-KEM-1024$^\dagger$ & $0{,}305$ & $0{,}007$ & $0{,}312$ & $2{,}2\%$ & $0{,}42$ & oui \\
SLH-DSA & $0{,}195$ & $0{,}006$ & $0{,}201$ & $3{,}0\%$ & $0{,}30$ & oui \\
\midrule
\multicolumn{7}{l}{\textbf{Sc\'enario bruit\'e} ($\eta_q = 0{,}15$)} \\
\midrule
ML-DSA-65 & $0{,}219$ & $0{,}012$ & $0{,}231$ & $5{,}2\%$ & $0{,}35$ & oui \\
ML-KEM-1024$^\dagger$ & $0{,}305$ & $0{,}014$ & $0{,}319$ & $4{,}4\%$ & $0{,}42$ & oui \\
\bottomrule
\end{tabular}
\end{table}

\noindent\textbf{Interpr\'etation du poids quantique.} Dans le sc\'enario
d'un support quantique ($\eta_q \approx 0{,}08$), la contribution du terme
quantique reste modeste, avec $2$ \`a $3\,\%$ de la majoration pour cette valeur contre $5\,\%$ pour
$\eta_q = 0{,}15$. Pour les signatures r\'eelles, $\eta_q = 0$ et la contribution
est nulle. La capacit\'e d'extraction reste le terme dominant.

\noindent\textbf{D\'epassement de la borne.} Les trois produits illustratifs
d\'epassent la majoration ($0{,}35 > 0{,}225$~; $0{,}42 > 0{,}312$~;
$0{,}30 > 0{,}201$). Ces syst\`emes ne sont
pas des instances du th\'eor\`eme (message public, $\Relver$ au lieu de $\Relrec$,
Remarque~\ref{rem:portee_cas_etude}) et les produits eux-m\^emes restent des
hypoth\`eses de travail. Le d\'epassement ne r\'efute donc pas le Trilemme. Il
emp\^eche d'affirmer que la borne est respect\'ee et d\'esigne $\Delta_q$ comme
premi\`ere cible du protocole de la section~\ref{subsec:protocole_delta_q}.

\end{sloppypar}

\begin{boiteArbitre}
L'arbitre s'arr\^ete sur un chiffre qui ne colle pas. Les produits sont
au-dessus de la borne, de dix \`a treize centi\`emes. Loin de le cacher, le texte
le nomme en donnant sa raison, puisque ces syst\`emes ne sont pas des instances
du th\'eor\`eme. L'arbitre note que la borne n'est ni confirm\'ee ni r\'efut\'ee
ici et qu'il reviendra \`a la mesure de $\Delta_q$ de trancher.
\end{boiteArbitre}

% -----------------------------------------------------------------------
\subsection{Propagation d'incertitude}
\label{sec:intervalles_bootstrap}
% -----------------------------------------------------------------------
\begin{sloppypar}

\subsubsection{M\'ethodologie}
\label{subsec:methodologie_bootstrap}

L'incertitude portant sur les param\`etres institutionnels est propag\'ee
jusqu'au rapport $\Delta/\lambda$ en affectant $\Delta$ d'un coefficient de
variation de $20\,\%$ pour les syst\`emes classiques et de $25\,\%$ pour les
syst\`emes post-quantiques, o\`u la p\'enalit\'e d'adaptation ajoute une source
d'incertitude. Ces coefficients sont des hypoth\`eses de travail~; ils ne sont
pas tir\'es d'une mesure.

\subsubsection{R\'esultats classiques}
\label{subsec:resultats_classiques_bootstrap}

Un coefficient de variation de $20\,\%$ sur $\Delta$ se propage int\'egralement
\`a $\Delta/\lambda$, dont l'intervalle \`a $95\,\%$ vaut environ $[0{,}17;\,0{,}38]$
(ECDSA), $[0{,}27;\,0{,}61]$ (Helios) et $[0{,}35;\,0{,}79]$ (zk-SNARKs). Cette
largeur traduit l'incertitude sur les param\`etres institutionnels~; elle
rappelle qu'en l'absence de donn\'ees, les valeurs de ce chapitre sont des
ordres de grandeur.

\begin{remark}[Intervalles de la formulation initiale]
\label{rem:ic_anterieurs}
La formulation initiale de cette section donnait des intervalles de demi-largeur
d'environ $1{,}8\,\%$ ($[0{,}268;\,0{,}278]$ pour ECDSA, par exemple), centr\'es sur
les valeurs calibr\'ees. Elle concluait que les \og mesures CLO\fg{} y tombaient.
Ces intervalles ne sont pas compatibles avec le coefficient de variation annonc\'e
(un tirage de $200\,000$ \'echantillons le montre, voir le script
\texttt{06\_calcul/python/verifications.py} du d\'ep\^ot)~; la comparaison portait
sur des valeurs identiques. La pr\'esente formulation ne les reprend pas.
\end{remark}

\subsubsection{R\'esultats quantiques}
\label{subsec:resultats_quantiques_bootstrap}

Avec un coefficient de variation de $25\,\%$ sur $\Delta_q$, le rapport
$\Delta_q/\lambda$ ($\lambda = 128$) a pour intervalle \`a $95\,\%$ environ
$[0{,}11;\,0{,}33]$ pour ML-DSA ($\Delta_q = 28$) et $[0{,}10;\,0{,}29]$ pour SLH-DSA
($\Delta_q = 25$). Le terme $\eta_q H(C)/\lambda$ est nul pour ces signatures,
qui sont des cha\^{\i}nes de bits classiques~; dans le sc\'enario hypoth\'etique
d'un support quantique ($\eta_q \approx 0{,}08$), il vaut moins d'un centi\`eme
et ne modifie pas ces intervalles. Les intervalles post-quantiques de
la formulation initiale ne sont pas repris, pour les m\^emes raisons.
\end{sloppypar}

% -----------------------------------------------------------------------
\subsection{Coh\'erence avec nos travaux ant\'erieurs}
\label{sec:validation_croisee_benchmarks}
% -----------------------------------------------------------------------
\begin{sloppypar}

\subsubsection{Rep\`eres classiques}
\label{subsec:benchmarks_classiques}

Le tableau~\ref{tab:benchmarks_classiques} situe les trois valeurs de $\Delta$
retenues par rapport aux rep\`eres propos\'es dans nos travaux ant\'erieurs et
dans la litt\'erature sur la rationalit\'e born\'ee. Hormis la r\'ef\'erence basse,
ces rep\`eres sont des auto-citations qui \'etablissent la coh\'erence avec nos
publications sans apporter de confirmation ind\'ependante.

\begin{table}[htbp]
\centering
\caption{Rep\`eres pour $\Delta$~: coh\'erence avec nos travaux ant\'erieurs}
\label{tab:benchmarks_classiques}
\begin{tabular}{@{}lcll@{}}
\toprule
\textbf{Institution} & $\Delta$ \textbf{estim\'e} &
\textbf{Source} & \textbf{Statut} \\
\midrule
Tribunaux petites cr\'eances & $\approx 25$ bits & \cite{Simon82} & Ordre de grandeur (analogie) \\
Autorit\'es fiscales & $\approx 65$ bits & \cite{MINKA2026CRO} & Travaux ant\'erieurs \\
R\'egulateurs financiers & $\approx 75$ bits & \cite{MINKA2026CRO} & Travaux ant\'erieurs \\
\midrule
ECDSA (tribunaux contrats) & 35 bits & \cite{MINKA2026SOTP} & Calibration \\
Helios (commissions \'elect.) & 56 bits & \cite{MINKA2026SOTP} & Calibration \\
zk-SNARKs (r\'egulateurs) & 73 bits & \cite{MINKA2026SOTP} & Calibration \\
\bottomrule
\end{tabular}
\end{table}

\subsubsection{Rep\`eres quantiques projet\'es}
\label{subsec:benchmarks_quantiques}

Aucun rep\`ere historique n'existe pour $\Delta_q$, car la transition PQC est en
cours et les institutions n'ont pas encore accumul\'e de pratique. Le
tableau~\ref{tab:benchmarks_quantiques} propose une grille de projection
index\'ee sur le niveau de pr\'eparation institutionnelle, qui rel\`eve de
l'hypoth\`ese.

\begin{table}[htbp]
\centering
\caption{Grille projet\'ee pour $\Delta_q$ selon le niveau de pr\'eparation PQC (hypoth\`ese)}
\label{tab:benchmarks_quantiques}
\begin{tabular}{@{}lcll@{}}
\toprule
\textbf{Niveau de pr\'eparation} & $\Delta_q$ \textbf{projet\'e} &
\textbf{Contexte} & \textbf{Horizon} \\
\midrule
Non pr\'epar\'e & $15$--$25$ bits & Transition d\'ebutante & 2025--2026 \\
Partiellement pr\'epar\'e & $25$--$45$ bits & Formation en cours & 2026--2027 \\
Bien pr\'epar\'e & $45$--$67$ bits & Standards int\'egr\'es & 2027--2028 \\
Expert PQC & $60$--$80$ bits & Anticipation avanc\'ee & 2028+ \\
\midrule
\multicolumn{4}{l}{\textit{Valeurs de calibration retenues}} \\
\midrule
ML-DSA-65 (non pr\'epar\'e) & 28 bits & Signature & 2025 \\
SLH-DSA (non pr\'epar\'e) & 25 bits & Signature (hachage) & 2025 \\
\bottomrule
\end{tabular}
\end{table}
\end{sloppypar}

% -----------------------------------------------------------------------
\subsection{\'Elicitation de plausibilit\'e}
\label{sec:evaluation_experts}
% -----------------------------------------------------------------------
\begin{sloppypar}

Les estimations ont fait l'objet d'\'echanges informels avec des interlocuteurs
de profils juridique, technique et r\'eglementaire. Faute d'un protocole
document\'e (fonctions des personnes consult\'ees, date, grille d'\'evaluation,
consentement), ces \'echanges ne sont pas pr\'esent\'es comme un r\'esultat de
cette th\`ese. Les scores de la formulation initiale ne sont pas repris. Une
\'elicitation structur\'ee, de type Delphi, pourrait \^etre int\'egr\'ee au
programme de la section~\ref{sec:protocoles_experimentaux}.
\end{sloppypar}

% -----------------------------------------------------------------------
\subsection{D\'ependance en $\lambda$}
\label{sec:stabilite_lambda}
% -----------------------------------------------------------------------
\begin{sloppypar}
\label{subsec:stabilite_classiques}\label{subsec:stabilite_quantiques}

\`A $\Delta$ fix\'e, la borne d\'ecro\^it en $1/\lambda$. Le rapport $\Delta/\lambda$
vaut ainsi $0{,}273$ pour ECDSA \`a $\lambda = 128$ et $0{,}137$ \`a $\lambda = 256$. Il s'agit
d'une pr\'ediction du mod\`ele, qui reste \`a tester. La formulation initiale
pr\'esentait des produits \og stables\fg{} en $\lambda$ ($0{,}273 \to 0{,}270$). La
forme de la borne conduit \`a ne pas reprendre ces tableaux. Que le produit CLO d'un
syst\`eme d\'eploy\'e varie peu avec $\lambda$ est plausible, mais cela illustre
justement que ces syst\`emes ne sont pas des instances du th\'eor\`eme
(Remarque~\ref{rem:portee_cas_etude}).
\end{sloppypar}

% -----------------------------------------------------------------------
\subsection{Test de significativit\'e}
\label{sec:test_significativite}
% -----------------------------------------------------------------------
\begin{sloppypar}
\label{subsec:test_classiques}\label{subsec:test_quantiques}

La formulation initiale rapportait une corr\'elation $r = 0{,}9997$ entre les deux
voies d'estimation, avec $p < 10^{-6}$ par permutation. Ce test n'est pas repris,
pour deux raisons. Avec trois points, un test de permutation ne compte que $3! = 6$
permutations, de sorte que la plus petite $p$-valeur atteignable est $1/6$. Surtout, la
corr\'elation mesurait une identit\'e, puisque $\Delta$ avait \'et\'e fix\'e par
$\Delta := \lambda \cdot (\Conf \cdot \Relver \cdot \Opp)$.

Pour les syst\`emes post-quantiques, la confrontation des produits CLO$^q$
illustratifs \`a la borne figure au tableau~\ref{tab:decomposition_borne_pq}, qui montre que
les trois produits la d\'epassent. Ce d\'epassement ne r\'efute pas le
th\'eor\`eme, dont ces syst\`emes ne sont pas des instances. Il emp\^eche en
revanche d'affirmer que la borne est \og respect\'ee\fg.
\end{sloppypar}

% -----------------------------------------------------------------------
\subsection{Donn\'ees brutes et format CSV}
\label{sec:donnees_brutes}
% -----------------------------------------------------------------------
\begin{sloppypar}

\subsubsection{Structure du jeu de donn\'ees}
\label{subsec:structure_csv}

Les valeurs de ce chapitre sont consign\'ees dans deux fichiers du d\'ep\^ot
compagnon. Le premier, \texttt{06\_calcul/data/parametres\_reference.csv}, reproduit le
tableau~\ref{tab:parametres_reference} et les triplets CLO des syst\`emes
classiques~; le second, \texttt{06\_calcul/data/calibration\_pq.csv}, reproduit le
tableau~\ref{tab:decomposition_borne_pq}. Leur sch\'ema complet figure dans
\texttt{06\_calcul/data/SCHEMA\_fr.md}~; le tableau~\ref{tab:structure_csv} en donne les
colonnes principales. Toutes les valeurs sont des param\`etres de calibration
illustratifs, sans mesure.

\begin{table}[htbp]
\centering
\caption{Colonnes principales des fichiers de donn\'ees}
\label{tab:structure_csv}
\begin{tabular}{@{}lll@{}}
\toprule
\textbf{Colonne} & \textbf{Type} & \textbf{Description} \\
\midrule
\texttt{system} & cha\^{\i}ne & Identifiant syst\`eme \\
\texttt{H\_C\_domaine}, \texttt{H\_C\_espece} & r\'eel & Entropie contextuelle (deux \'echelles) \\
\texttt{Delta}, \texttt{Delta\_q} & r\'eel & Capacit\'e d'extraction (calibration) \\
\texttt{alpha}, \texttt{beta} & r\'eel & Granularit\'e, r\'egularit\'e \\
\texttt{eta\_q} & r\'eel & Perte d'interpr\'etabilit\'e (sc\'enario~; $0$ pour une cha\^{\i}ne de bits) \\
\texttt{lambda} & entier & Param\`etre de s\'ecurit\'e \\
\texttt{product\_clo} & r\'eel & Produit CLO (calibration illustrative, $\Relver$) \\
\texttt{borne\_these} & r\'eel & Majoration affaiblie $(\Delta_q + \eta_q H(C))/\lambda$ \\
\bottomrule
\end{tabular}
\end{table}

\label{subsec:echantillon_classiques}\label{subsec:echantillon_quantiques}
Les colonnes \texttt{product\_pred}, \texttt{product\_meas} et \texttt{borne} de la
formulation initiale deviennent \texttt{product\_clo} et \texttt{borne\_these}, car
le jeu de donn\'ees ne contient aucune mesure et la \og pr\'ediction\fg{} co\"{\i}ncidait
avec le produit CLO.
\end{sloppypar}

% -----------------------------------------------------------------------
\subsection{Synth\`ese}
\label{sec:synthese_validation}
% -----------------------------------------------------------------------
\begin{sloppypar}

Le tableau~\ref{tab:synthese_validation} r\'ecapitule ce que cette section
\'etablit ainsi que ce qu'elle n'\'etablit pas.

\begin{table}[htbp]
\centering
\caption{Synth\`ese de la calibration r\'etrospective}
\label{tab:synthese_validation}
\small
\begin{tabular}{@{}llccl@{}}
\toprule
\textbf{Crit\`ere} & \textbf{M\'etrique} & \textbf{Classique} & \textbf{PQC} &
\textbf{Statut} \\
\midrule
Plausibilit\'e des $\Delta$ & Intervalle cognitif $[20,80]$ & 3/3 & 2/2 & Ordre de grandeur \\
Hypoth\`ese~2$'$ & $H(V|C) \geq \gamma$ & 1/3 (ECDSA) & --- & Hors champ sinon \\
Produit $\leq$ borne & --- & non applicable & d\'epass\'ee & Syst\`emes hors champ \\
Incertitude & IC~95\,\% de $\Delta/\lambda$ & $\pm 39\,\%$ & $\pm 49\,\%$ & CV suppos\'e \\
\bottomrule
\end{tabular}
\end{table}

\noindent\textbf{Conclusion.}
\begin{itemize}
\item \textbf{Classique (C2/C4).} Les ordres de grandeur retenus sont
  mutuellement coh\'erents~; aucune validation empirique n'est revendiqu\'ee.
\item \textbf{Post-quantique (C3).} Les produits illustratifs d\'epassent la
  borne~; les param\`etres $\Delta_q$ sont des postulats \`a mesurer en priorit\'e.
\item \textbf{R\'eserve.} La validation prospective demeure n\'ecessaire. Les
  protocoles de la section~\ref{sec:protocoles_experimentaux} fournissent les
  mesures de $\Delta$, $\Delta_q$ et $\eta_q$ qui restent \`a produire, avec les
  crit\`eres de r\'efutation du tableau~\ref{tab:criteres_refutation}~; leur
  statut de d\'eploiement est \'etabli \`a la
  section~\ref{subsubsec:statut_deploiement}.
\end{itemize}
\end{sloppypar}

% -----------------------------------------------------------------------
\subsection{Reproductibilit\'e}
\label{sec:reproductibilite}
% -----------------------------------------------------------------------
\begin{sloppypar}

\subsubsection{D\'ep\^ot}
\label{subsec:depot}

\begin{center}
\url{https://github.com/MINKAThierry/PhDThesis}\\
(version 1.0.0, archiv\'ee sous le DOI \texttt{10.5281/zenodo.XXXXXXX})
\end{center}

Le d\'ep\^ot~\cite{CROTrilemmaDepot} s'organise en huit r\'epertoires, d\'ecrits au
tableau~\ref{tab:depot_structure}. Chacun contient un fichier \texttt{README} bilingue
(fran\c{c}ais et anglais) qui indique le statut de son contenu, publi\'e, pr\'eprint ou
proposition de th\`ese.

\begin{table}[htbp]
\centering
\caption{Organisation du d\'ep\^ot}
\label{tab:depot_structure}
\small
\begin{tabularx}{\textwidth}{@{}l X l@{}}
\toprule
\textbf{R\'epertoire} & \textbf{Contenu} & \textbf{Partie de la th\`ese} \\
\midrule
00\_these & Sources \LaTeX{} et PDF du manuscrit & Ensemble \\
01\_publie & Les trois articles publi\'es dans IEEE Access (P1, P2, P3) & Chap.~1 \`a 3, annexe~\ref{chap:annexe_publications} \\
02\_preprints & Fiches des neuf pr\'eprints IACR ePrint cit\'es & Chap.~1, annexe~\ref{chap:annexe_publications} \\
03\_resultats\_these & \'Enonc\'es propres \`a la th\`ese et formulations pr\'ecis\'ees par rapport aux articles & Chap.~2 \`a 4 \\
04\_sotp & Sp\'ecification de la famille SOTP, impl\'ementations de r\'ef\'erence et preuves game-based & Chap.~\ref{chap:constructions}, annexe~\ref{chap:annexe_gamebased} \\
05\_formalisation\_lean & Fichier ConditionalCRO.lean (version 2.2) et configuration Lean/Mathlib & Section~\ref{sec:lean4}, annexe~\ref{chap:annexe_lean4} \\
06\_calcul & Scripts, donn\'ees de calibration, sorties de r\'ef\'erence et tests & Section~\ref{sec:validation_retrospective} \\
07\_protocoles\_terrain & Instruments des protocoles P1 \`a P7, non d\'eploy\'es & Section~\ref{sec:protocoles_experimentaux} \\
\bottomrule
\end{tabularx}
\end{table}

Le r\'epertoire \texttt{06\_calcul/} r\'eunit les
scripts ci-dessous (\texttt{python/}), les donn\'ees de calibration (\texttt{data/}),
les sorties de r\'ef\'erence avec leurs empreintes SHA-256 (\texttt{results/}) et une
batterie de tests qui recalcule chaque valeur num\'erique cit\'ee dans ce chapitre
(\texttt{tests/}). La commande \texttt{./reproduire.sh}, lanc\'ee dans
\texttt{06\_calcul/}, r\'eg\'en\`ere l'ensemble.

\subsubsection{Environnement}
\label{subsec:environnement}

\begin{table}[htbp]
\centering
\caption{Environnement de calcul pour reproductibilit\'e}
\label{tab:environnement_calcul}
\begin{tabular}{@{}ll@{}}
\toprule
\textbf{Composante} & \textbf{Version} \\
\midrule
Python & 3.10 \`a 3.12 (r\'ef\'erence~: 3.11) \\
NumPy & $\geq 1.24$ \\
mpmath & $\geq 1.3$ \\
Lean~4 / Mathlib & v4.34.0 / v4.34.0 \\
\bottomrule
\end{tabular}
\end{table}

\subsubsection{Scripts}
\label{subsec:scripts_validation}

Les commandes suivantes s'ex\'ecutent depuis \texttt{06\_calcul/python/}.

\noindent\textbf{Calibration classique~:}
\begin{verbatim}
$ python validate_sotp.py --system all --type classical --lambda 128
\end{verbatim}

\noindent\textbf{Calibration post-quantique~:}
\begin{verbatim}
$ python validate_sotp.py --system all --type postquantum --lambda 128
\end{verbatim}

\noindent\textbf{Analyse de sensibilit\'e $\eta_q$ (sc\'enario hypoth\'etique)~:}
\begin{verbatim}
$ python sensitivity_eta_q.py --range 0.01:0.20 --steps 20
\end{verbatim}
\end{sloppypar}

\begin{boiteArbitre}
L'arbitre pose la plume. Il a vu des ordres de grandeur coh\'erents entre eux,
des intervalles dont la largeur est assum\'ee sans fard ainsi que des d\'epassements de borne
nomm\'es comme tels. Il a vu aussi des tests \'ecart\'es parce qu'ils
n'apportaient aucune information.

Il retient surtout la derni\`ere ligne, selon laquelle la validation prospective reste
\`a faire. Elle est nomm\'ee comme telle, avec des crit\`eres de r\'efutation
fix\'es \`a l'avance. L'arbitre ne confond pas calibration et preuve~; le dossier non
plus.
\end{boiteArbitre}
